\documentclass[12pt]{article}

\usepackage[a4paper, total={6in, 8in}]{geometry}
\usepackage{textcomp}

\begin{document}
\setlength{\parindent}{0pt}

\def \t {\textrightarrow}
\def \v {\vspace{1em}}
\def \bi {\begin{itemize}}
\def \ei {\end{itemize}}
\def \s[#1] {\section*{#1}}
\def \ss[#1] {\subsection*{#1}}
\def \sss[#1] {\subsubsection*{#1}}

\s[Appunti 21 marzo su foglio a penna]

\s[Architettura fascista]
\ss[Casa del fascio - Terragni]
Capolavoro del razionalismo fascista, a Como \\
Terragni è lo stesso della casa Rustici a Milano \t studia a como ma si laurea al politecnico di milano \\
Sara sulla nave della Carta d'Atene \t scrive riviste, e lavora soprattutto a Como (questa opera, poi fa Novum Covum, asilo di Sant'Elia) \\
Poi va in guerra (seconda guerra mondiale) \t andrà anche in Russia \t muore giovane nel 43 anni per trombosi cerebrale \\
Diventa la sede del partito fascista \t ora è sede della guardia di finanza \\
Dal punto di vista funzionale ci sono quindi gli uffici e la sala delle adunanze

\v

Al piano terra ci sono corridoi, a primo piano ballatoi che si affacciano sulla sala delle adunanze \t e ballatoi collegano i diversi uffici \\
Non c'è un trattamento monumentale di una facciata rispetto alle altre \t no facciata principale vera \\
4 piani, che si organizza intorno a una corte decentrata \t la corte è la sala delle adunanza \\
Collegamento verticale è una scala \\
Struttura a corte richiama alla tradizione architettonica italiana \t tipica del palazzo rinascimentale \\
La progettazione si basa sulla geometria \t usa il quadrato come forma base \t rimanda a rinascimento con Alberti \t e il palazzo è un mezzo cubo

\ss[Struttura]
Punto di vista strutturale ci sono pilastri e travi in calcestruzzo armato \t e sono a vista, sono visibili + tamponamenti in laterizio \\
La copertura della sala delle adunanza è il vetro cemento \t brevettato a inizio 900, usato principalmente in edifici industriali \\
Il vetro cemento \t sono due formelle di vetro, con intercapedine di aria \t e poi vengono unite con cemento \t proprietà è che fa passare la luce, ma impedisce di vedere bene all'interno \t ed è molto resistente \\
Materiale nuovo che usa \\
Il rivestimento esterno è tutto in marmo bianco \t il movimento moderno non usa materiali preziosi come il marmo \t anche in una struttura cosi geometrica ed essenziale da + monumentalita, ma lui giustifica dicendo che marmo è molto resistente \

\ss[Interno]
L'interno è intonacato in verde pastello \t ora invece è bianca \\
La scala è molto essenziale, il sistema strutturale è esibito \t nessun tipo di decorazione, ne interno ne esterno

\v

I prospetti sono tutti diversi \t no facciata principale \t sono tutti con marmo e studiati su base geometrica \\
Suddivide geometricamente dei rettangoli, alti 4 e lunghi 7 \t lascia i due rettangoli laterali coperti, mentre negli altri ci sono logge con struttua portante a vista \\
Progettazione su base geometrica evidente

\s[Novecento]
Novecento è un gruppo di artisti e architetti che si raccoglie attorno alla figura di Margherita Sarfatti \t a Milano, lei è una critica d'arte \\
Viene da famiglia venenziana ebrea? \t sara per molti anni amante di mussolini, che conosce alla redazione dell'Avanti \\
Relazione finira, e lei andra via dall'italia con le leggi razziali

\v

Questi artisti si pongono l'obbiettivo di ricreare un arte italiana \t in realta c'è fenomeno degli anni 20 chiamato "Ritorno all'ordine" \\
Dopo 1 guerra mondiale europa è distrutta psicologicamente \t c'è bisogno di certezze e serenità, che non si trova nelle avanguardie, che esprimono la crisi delle certezze di primo novecento \t non rasserenano \\
Si vede quindi ritorno al naturalismo, alla prospettiva etc.

\ss[L'allieva - Sironi]
Sironi è futurista, è uno dei protagonisti \t quest'opera è del 24 \t ritratto con figura femminile, statica \\
Forte chiaro/scuro \t e dietro c'è una statuetta che è un nudo che rimanda alla classicita \\
Sironi sara anche principale esponente dell'arte figurativa fascista

\ss[Donne che corrono sulla spiaggia - Picasso]
Lui ha percorso artistico molto lungo che comprende molte fasi \t reinterpreta novita, etc. \\
In questo periodo ha periodo neoclassico \t anche lui partecipa in questo ritorno all'ordine \\
La destrutturazione cubista non è presente \t sono due figure femminili (sproporzionate e deformate, ma lo stesso che rimandano a un naturalismo)

\v

Si vuole creare quindi un arte fascita e italiana, e ci si rifa quindi al rinascimento e alla classicita \t ma non vuol dire imitazione, solo riferimento \\
Dal punto di vista architettonico \t stesso discorso: si riferisce alla romanita e alla classicita, ma non c'è imitazione (no neoclassicismo, ma un utilizzo degli elementi del linguaggio classico che vengono pero semplificati) \\
Abbiamo Portalupi, Giò Ponti, etc. \t fondata rivista "Domus", diretta da Giò Ponti

\sss[Palazzo dell'Arte - Giovanni Muzio]
Ha una funzione espostiva, e viene realizzato da Muzio quando l'esposizione triennale nel 33 viene trasferita a Milano (prima Monza) \\
Palazzo espositivo, che utilizza il linguaggio Novecento \t parte posteriore geometrica ed essenziale, con porticato a forma di C verso il giardino \\
È una successione di arcata a tutto sesto \t che rimanda alla romanita, ma è estremamente essenziale \t non ci sono capitelli, colonne etc.
\end{document}
